\begin{table}[H]
\small
\setlength{\tabcolsep}{4.2pt}
\renewcommand{\arraystretch}{0.95}

\caption{\label{tab:ea20-income-gap-product-contributions}Product-group contributions to euro-area cumulative income inflation gaps, January 2021--December 2023}
\centering
\resizebox{\textwidth}{!}{%
\begin{tabular}{lrrrrr}
\toprule
Product group & Product weight & Cumulative inflation & Q1 contribution & Q5 contribution & Inflation gap\\
\midrule
Food and non-alcoholic beverages & 16.6 & 4.4 & 5.2 & 3.8 & 1.4\\
Alcoholic beverage, tobacco and narcotics & 4.3 & 0.6 & 0.9 & 0.4 & 0.4\\
Clothing and footwear & 5.2 & 1.0 & 0.8 & 1.1 & -0.3\\
Housing (rentals and repairs) & 12.2 & 1.7 & 1.4 & 1.8 & -0.4\\
Water, electricity, gas and other fuels & 6.3 & 3.4 & 4.6 & 2.7 & 1.9\\
\addlinespace
Fuel & 5.9 & 1.8 & 1.4 & 1.9 & -0.5\\
Transport (excluding fuel) & 8.9 & 0.8 & 0.5 & 1.1 & -0.5\\
Other & 40.5 & 4.1 & 3.3 & 4.6 & -1.4\\
\midrule
\textbf{Total} & \textbf{100.0} & \textbf{17.8} & \textbf{18.2} & \textbf{17.5} & \textbf{0.6}\\
\bottomrule
\end{tabular}%
}
\vspace{-0.4em}
\begin{flushleft}
\footnotesize{Notes. Product contributions are derived from the same observed product-country indices as Table~1. Product indices are aggregated within each country before national monthly movements are aggregated to EA20 with annual country weights. Cumulative changes compare January 2021 with December 2023. The chain-linked decomposition is exact; no normalization or residual allocation is used. The gap is Q1 minus Q5. Weights are 2022 EA20 expenditure shares. Sources: Eurostat, authors' calculations.}
\end{flushleft}
\end{table}
